% ch10_b §10.4–10.5 (书页 336–340 / p350–p354)
%--A-TAIL--
%FIG{176}: 铁磁体中的自发磁化。

这时系统的行为就好像每一个磁矩都平行排列一样；我们说这固体是\emph{铁磁性的}（ferromagnetic），并且看到，即使在不存在磁场时它也表现出很大的自发磁化。当我们使晶体升温时，自发磁化不断减小，直到在一个完全确定的温度——\emph{居里温度}（Curie temperature）——处消失：
\begin{equation*}
T_{\mathrm{C}}=\theta.
\tag{10.28}
\end{equation*}

\section{交换相互作用}

用居里温度 (10.25) 来衡量，Weiss 内场非常之大——远大于一个磁体阵列的普通内禀偶极场。在许多铁磁物质中，系数 $\lambda$ 并非 $\frac{4}{3}\pi$（对 Lorentz 场 (8.52) 而言它本可以取这个值），而是 1000 的量级。

通常把这个表观场归因于相邻原子上电子自旋之间的\emph{交换相互作用}（exchange interaction）。对于 Heisenberg 和 Dirac 的严格局域模型，我们可以给第 $\boldsymbol{l}$ 个格点上的原子指派一个总自旋算符 $\boldsymbol{\mathsf{S}}_{\boldsymbol{l}}$，并写下哈密顿量
\begin{equation*}
\mathcal{H}=-\sum_{\boldsymbol{l}\boldsymbol{l}'}J_{\boldsymbol{l}\boldsymbol{l}'}\,\boldsymbol{\mathsf{S}}_{\boldsymbol{l}}\cdot\boldsymbol{\mathsf{S}}_{\boldsymbol{l}'}-\beta\mathbf{H}\cdot\sum_{\boldsymbol{l}}\boldsymbol{\mathsf{S}}_{\boldsymbol{l}}.
\tag{10.29}
\end{equation*}
\emph{交换积分}（exchange integral）$J_{\boldsymbol{l}\boldsymbol{l}'}$ 是格点 $\boldsymbol{l}$ 与 $\boldsymbol{l}'$ 相对位置的函数，但只有当 $\boldsymbol{l}-\boldsymbol{l}'$ 为一两个格点间距时它才比较大。

我们可以把这个哈密顿量写成
\begin{equation*}
\mathcal{H}=-\sum_{\boldsymbol{l}}\bigl\{\sum_{\boldsymbol{l}'}J_{\boldsymbol{l}\boldsymbol{l}'}\boldsymbol{\mathsf{S}}_{\boldsymbol{l}'}+\beta\mathbf{H}\bigr\}\cdot\boldsymbol{\mathsf{S}}_{\boldsymbol{l}}.
\tag{10.30}
\end{equation*}
试想温度极低时会出现什么情形：那时所有自旋都排列整齐，以致对一切 $\boldsymbol{l}'$ 有
\begin{equation*}
\langle\beta\boldsymbol{\mathsf{S}}_{\boldsymbol{l}'}\rangle\approx\langle\boldsymbol{\mu}\rangle,
\tag{10.31}
\end{equation*}
于是我们可以把 (10.30) 中的第一项辨认成内场
\begin{align*}
\beta\mathbf{H}_{I}&=\sum_{\boldsymbol{l}'}J_{\boldsymbol{l}\boldsymbol{l}'}\boldsymbol{\mathsf{S}}_{\boldsymbol{l}'}\\
&=\frac{1}{\beta}\Bigl(\sum_{\boldsymbol{l}'}J_{\boldsymbol{l}\boldsymbol{l}'}\Bigr)\langle\boldsymbol{\mu}\rangle,
\tag{10.32}
\end{align*}
即
\begin{equation*}
\lambda=\frac{1}{N\beta^{2}}\Bigl(\sum_{\boldsymbol{l}'}J_{\boldsymbol{l}\boldsymbol{l}'}\Bigr)
\tag{10.33}
\end{equation*}
或者，用居里温度 (10.25) 来表示，
\begin{equation*}
k\theta=\frac{1}{3}S(S+1)\sum_{\boldsymbol{l}'}J_{\boldsymbol{l}\boldsymbol{l}'}.
\tag{10.34}
\end{equation*}
这便用系统转变为铁磁态所处的温度给出了交换相互作用的大小。

这个\emph{自旋哈密顿量}（spin Hamiltonian）(10.29) 的性质将是本章其余部分的主要论题。但是，它作为一个基本假设是否有充分的根据呢？对这个问题已有过许多不同的回答，而它正是铁磁性作为一种物理现象之谜的核心所在。

从群论的论证容易证明：相互作用取自旋算符的标量积这一形式，至少是最简单的可能形式，尽管还可能存在更复杂的其他项，例如\emph{偶极–偶极相互作用}（dipole–dipole interaction），
\begin{equation*}
\mathcal{H}_{\mathrm{dip.-dip.}}=\beta^{2}\biggl\{\frac{\boldsymbol{\mathsf{S}}_{1}\cdot\boldsymbol{\mathsf{S}}_{2}}{R^{3}}-\frac{3(\boldsymbol{\mathsf{S}}_{1}\cdot\mathbf{R})(\boldsymbol{\mathsf{S}}_{2}\cdot\mathbf{R})}{R^{5}}\biggr\},
\tag{10.35}
\end{equation*}
其中两个自旋相距 $\mathbf{R}$（一个矢量）。

同样众所周知的是，同一原子上或相邻原子上的电子自旋倾向于通过交换效应耦合起来——这是 Pauli 原理的一个推论。例如，若 $\phi_{a}$ 和 $\phi_{b}$ 是我们“放入两个电子”的两个波函数，那么按照这两个电子的自旋是平行还是反平行，我们可以构造出两类状态。这些状态就是对称与反对称组合
\begin{align*}
\Phi_{S}(\mathbf{r}_{1},\mathbf{r}_{2})&=2^{-\frac{1}{2}}\bigl\{\phi_{a}(\mathbf{r}_{1})\phi_{b}(\mathbf{r}_{2})+\phi_{a}(\mathbf{r}_{2})\phi_{b}(\mathbf{r}_{1})\bigr\},\\
\Phi_{A}(\mathbf{r}_{1},\mathbf{r}_{2})&=2^{-\frac{1}{2}}\bigl\{\phi_{a}(\mathbf{r}_{1})\phi_{b}(\mathbf{r}_{2})-\phi_{a}(\mathbf{r}_{2})\phi_{b}(\mathbf{r}_{1})\bigr\},
\tag{10.36}
\end{align*}
其中 $\Phi_{S}$ 属于反平行自旋（\emph{单态}，singlet state），而 $\Phi_{A}$ 属于平行自旋（\emph{三重态}，triplet state）。现在，如果我们计算库仑能 $e^{2}/|\mathbf{r}_{1}-\mathbf{r}_{2}|$ 在这两个态中的平均值，就会发现二者相差一个量
\begin{equation*}
2\iint\phi_{a}^{*}(\mathbf{r}_{1})\phi_{b}^{*}(\mathbf{r}_{2})\frac{e^{2}}{|\mathbf{r}_{1}-\mathbf{r}_{2}|}\phi_{a}(\mathbf{r}_{2})\phi_{b}(\mathbf{r}_{1})\,\mathrm{d}\mathbf{r}_{1}\,\mathrm{d}\mathbf{r}_{2},
\tag{10.37}
\end{equation*}
这就是交换积分。容易把这个差改写成如下形式，使它表示 $2J\,\boldsymbol{\mathsf{S}}_{1}\cdot\boldsymbol{\mathsf{S}}_{2}$ 在 $\boldsymbol{\mathsf{S}}_{1}$ 与 $\boldsymbol{\mathsf{S}}_{2}$ 平行和反平行两种情形下取值之差。于是，$J$ 的正负与大小便取决于这个积分的正负与大小。

有两种众所周知的情形。设两个电子属于同一原子，但并未填满一个闭合壳层。这时，库仑相互作用与原子轨道 $\phi_{a}$、$\phi_{b}$ 的形式使得 $J$ 为正；电子自旋倾向于排列起来，在与壳层中待填充的独立态数目相容的前提下使总自旋最大。这就是 Hund 规则（Hund's rule），它解释了为什么过渡金属离子的未满 d 壳层中的电子倾向于组合起来，赋予离子一个很大的永久磁矩（§5.5）。研究这类离子处于复杂晶体中各种环境时的顺磁性与磁共振性质，是物理学的一个重要分支。

但我们在这里关心的是不同离子上的电子自旋之间的相互作用——结果发现，在这种情形下算出的 $J$ 几乎总是负的，即倾向于使相邻格点上的自旋反平行。这方面最简单的例子是氢分子的 Heitler–London 模型，其中成键态的电子自旋是配对的。在 §4.2 中处理共价键时，我们曾把这一点视为理所当然。

因此，在这个基础上很难解释许多金属——通常是“过渡”族元素——具有铁磁性这一事实。我们可以理解为什么每个离子中的 d 电子倾向于组合成一个总自旋——在 Fe 中大到 $\frac{5}{2}$；却无法理解为什么相邻离子上的磁矩之间的相互作用会有利于平行取向。问题还因如下事实（§4.1）而更加复杂：这些所谓的 d 电子并非严格局域在特定的离子上，而是处于从原子到原子相互交叠、形成窄能带的状态。这个窄带又与普通的 s 带杂化，在 s 带中电子可以十分自由地传导。Heisenberg 模型哈密顿量 (10.29) 尽管作为金属铁磁性的唯象理论颇为成功，却需要从第一性原理出发加以认真的重新考察。这个微妙问题——固体理论中最重要的基本问题之一——的若干方面将在 §10.5–§10.6 中讨论。

\section{能带铁磁性}

从 Heisenberg 模型的局域自旋转向相反的极端，让我们把所有磁性电子都放进一个由 Bloch 态构成的能带。令 $n_{\mathbf{k}+}$ 为具有 + 自旋的态 $|\mathbf{k}\rangle$ 的占据数。在 Stoner 的\emph{集体电子模型}（collective electron model）中，我们在电子哈密顿量上添加如下一项：
\begin{equation*}
\mathcal{H}_{\mathrm{int.}}=\frac{U}{N}\sum_{\mathbf{k},\mathbf{k}'}n_{\mathbf{k}+}n_{\mathbf{k}'-}.
\tag{10.38}
\end{equation*}
于是，每一对自旋相反的电子都贡献一个大小为 $U/N$ 的正的“交换”能，它或许来自这两个电子偶尔同时栖身于同一原子的 d 壳层，但不一定恰好由像 (10.37) 那样的积分严格表示。自旋相同的电子的贡献则无须显式计入，因为它们可以归并到能量零点的定义之中。

虽然这一项可以解释为给出一个像 (10.22) 那样的内场，但我们现在处理的是 费米–狄拉克 分布，可以沿用 §10.2 的论证。令 $n_{\sigma}$ 为每原子中自旋为 $\sigma$ 的电子总数。代替 (10.9)，一个 + 自旋电子的能量为
\begin{equation*}
\mathcal{E}_{\mathbf{k}+}=\mathcal{E}(\mathbf{k})-\mu_{0}H+Un_{-},
\tag{10.39}
\end{equation*}
对 $-$ 自旋的电子有类似的式子。

两个电子集合的化学势仍然相同，因此我们必须满足如下方程：
\begin{equation*}
n\langle\mu\rangle=\mu_{0}\int_{0}^{\infty}\frac{1}{2}\bigl\{f^{0}(\mathcal{E}_{\mathbf{k}+})-f^{0}(\mathcal{E}_{\mathbf{k}-})\bigr\}\mathcal{N}(\mathcal{E})\,d\mathcal{E},
\tag{10.40}
\end{equation*}
以及
\begin{equation*}
n=\int\frac{1}{2}\bigl\{f^{0}(\mathcal{E}_{\mathbf{k}+})+f^{0}(\mathcal{E}_{\mathbf{k}-})\bigr\}\mathcal{N}(\mathcal{E})\,d\mathcal{E}.
\tag{10.41}
\end{equation*}
在 $H\to 0$、$T\to 0$ 的极限下，这些方程有一个与 (10.11) 类似的解。磁化率取如下形式：
\begin{equation*}
\chi=\frac{\mu_{0}^{2}\mathcal{N}(\mathcal{E}_{F})}{1-\frac{1}{2}U\mathcal{N}(\mathcal{E}_{F})}
\tag{10.42}
\end{equation*}
就好像我们只是把 Pauli 磁化率 (10.12) 增大了一个因子 $1/\{1-\frac{1}{2}U\mathcal{N}(\mathcal{E}_{F})\}$。交换相互作用 (10.38) 由于偏向平行自旋，使系统更容易被极化。

%FIG{177}: (a) 能带铁磁性；(b) 铁磁金属的电子比热。

但是，当交换场强到足以使
\begin{equation*}
\frac{1}{2}U\mathcal{N}(\mathcal{E}_{F})>1,
\tag{10.43}
\end{equation*}
时，这个形式解显然是不稳定的。这标志着向铁磁性的转变。为使 (10.39)–(10.41) 相互一致，我们必须让 $n_{+}$（比如说）远大于 $n_{-}$，从而使系统具有一个永久磁矩。随后便可数值求解这些方程，给出 $\langle\mu\rangle$ 作为温度的函数。其结果在定性上与 (10.26) 的结果并无明显不同。系统表现出一个居里温度，在居里温度以下磁矩饱和到近乎恒定的值；在居里点以上，理论也预言了某种类似于居里–外斯定律 (10.24) 的行为。

在这个模型中，不难显式计算出与这一转变相联系的比热。能量（参见 §4.7）可以作为 $T$ 的函数算出，再对它求导便给出比热。结果表明，$C_{\mathrm{el.}}$ 在居里